\chapter{Introduction}
This chapter serves as an introduction into the master's thesis, explaining why 
the topic of sleep stage classification is important and interesting.
We describe the gap in existing research that we aim to fill
and quickly describe the used approach to tackle the problem.

\section{Motivation}
Sleep staging is the process of analyzing data recorded during sleep to classify
different stages of sleep. This classification is used in both clinical and research
settings to understand sleep patterns, diagnose sleep disorders, and evaluate the
effectiveness of treatments. Traditional sleep staging is performed manually by trained
professionals, who have to meticulously check and mark the recordings. Reportedly, thousands of hours of sleep are recorded nightly \cite{VallatWalker2021}. Additionally,
on average, professionals agree on only about 83\% of the sleep stages when scoring the same
recordings \cite{Rosenberg2013}. This makes the process both time-consuming and subjective.

In recent years, various machine learning solutions have been implemented 
to help sleep technicians make sleep staging more efficient and less burdensome.
While accuracy of these solutions is usually high, they are often designed 
to solve a specific subproblem, such as classifying patients in particular age groups.
However, sleep patterns can vary significantly across different demographics. For example,
total power of lower frequency brain waves, which are primarily used to classify the sleep stage,
has been shown to decrease with age \cite{Greve2025}.

Furthermore, the transformer architecture \cite{Vaswani2017AttentionIA} which
has shown great success in other temporal data domains, remains a
relatively unexplored approach for sleep staging.

This master's thesis aims to develop and evaluate a generalizable sleep staging model
that can accurately classify sleep stages across diverse demographics. We do this by proposing multiple
strategies, each based on the transformer architecture, then compare and discuss their effectiveness.

